% !TeX root = ../../Book2.tex
% 纲要标题：### C.2 Azumaya 代数的局部性质
% 纲要位置：工作记录/计划纲要.md，第 1834 行。

% 纲要要点（供目录核对）：
% * 局部化与剩余域纤维分别检验；Nakayama、目标投射性和迹单位的使用位置，主开有限覆盖、任意变基、中心、泛纤维反例与单位判据；平展正向提升明确外部归属，反向完整下降
% #### C.2.1 有限投射模的局部检验
% #### C.2.2 纤维判据与底环变换

\section{Azumaya 代数的局部性质}
\label{b2:sec:C02}

当代数的底模已知忠实、有限生成且投射时，左右乘作用是两个有限投射模之间的映射，因而可以在剩余域上检验。本节建立这一检验所需的局部模论。这里的局部化使用\cref{b2:def:localized-module}的分式构造，零判据为式\eqref{b2:eq:localization-zero}；局部化保持正合性，来自\cref{b2:prop:localization-flat}与其张量描述。

\subsection{有限投射模的局部检验}
\label{b2:subsec:C02:01}

设 \(R\) 为非零交换含幺环。\(R\) 的全部素理想组成的集合称为\term[sulixiang-pu]{素理想谱}[prime spectrum]，记为 \(\operatorname{Spec}R\)。对 \(\mathfrak p\in\operatorname{Spec}R\)，令 \(R_{\mathfrak p}=(R\setminus\mathfrak p)^{-1}R\)，并对 \(R\) 模 \(M\) 令 \(M_{\mathfrak p}=R_{\mathfrak p}\otimes_R M\)。环 \(R_{\mathfrak p}\) 的唯一极大理想为 \(\mathfrak pR_{\mathfrak p}\)：分子不在 \(\mathfrak p\) 中的分式可逆，分子在 \(\mathfrak p\) 中的分式都属于这个真理想。它的\term[shengyu-yu]{剩余域}[residue field]及 \(M\) 在该点的\term[xianwei-mo]{纤维}[fiber of a module]分别记为
\[
\kappa(\mathfrak p)=R_{\mathfrak p}/\mathfrak pR_{\mathfrak p},
\qquad M(\mathfrak p)=\kappa(\mathfrak p)\otimes_R M.
\]
若 \(\mathfrak m\) 极大，则 \(\kappa(\mathfrak m)=R/\mathfrak m\)。这里把素理想看作一个点，剩余域是在该点取值所用的系数域，而纤维是这个域上的向量空间。局部化保留了局部环上的模，取纤维还要模去其极大理想：
\[
M(\mathfrak p)\cong M_{\mathfrak p}/\mathfrak pM_{\mathfrak p}.
\]
\symidx{\(\operatorname{Spec}R\)：交换环的素理想谱}
\symidx{\(\kappa(\mathfrak p)\)：素理想处的剩余域}
\symidx{\(M(\mathfrak p)\)：模在素理想处的纤维}

\begin{lemma}[局部消失与有限投射模]\label{b2:lem:azumaya-local-toolkit}
设 \(R\) 为非零交换含幺环。对 \(R\) 模 \(X\) 和素理想 \(\mathfrak p\)，记 \(R_{\mathfrak p}=(R\setminus\mathfrak p)^{-1}R\)、\(X_{\mathfrak p}=R_{\mathfrak p}\otimes_RX\)、\(\kappa(\mathfrak p)=R_{\mathfrak p}/\mathfrak pR_{\mathfrak p}\)，并记纤维 \(X(\mathfrak p)=\kappa(\mathfrak p)\otimes_RX\)。
\begin{enumerate}
\item 对任意 \(R\) 模 \(M\)，\(M=0\) 当且仅当每个极大理想 \(\mathfrak m\) 处的 \(M_{\mathfrak m}=0\)。若 \(M\) 有限生成，还可等价地要求全部纤维 \(M(\mathfrak m)=0\)。
\item 对任意有限生成投射 \(R\) 模 \(P\)，每个 \(P_{\mathfrak p}\) 都是有限自由 \(R_{\mathfrak p}\) 模。
\item 若 \(P\) 为有限生成 \(R\) 模、\(Q\) 为有限生成投射 \(R\) 模，且 \(u:P\rightarrow Q\) 为 \(R\) 模同态，则 \(u\) 为同构，当且仅当它在每个极大理想 \(\mathfrak m\) 处诱导的 \(u(\mathfrak m):P(\mathfrak m)\rightarrow Q(\mathfrak m)\) 为同构。
\end{enumerate}
\end{lemma}
\begin{proof}
若 \(0\ne x\in M\)，要找一个局部化仍能看见它的点，只需让分母避开它的零化子。具体地，\(I=\{r\in R:rx=0\}\) 是真理想，取包含它的极大理想 \(\mathfrak m\)。若 \(x/1=0\) 于 \(M_{\mathfrak m}\)，分式零判据给出 \(s\notin\mathfrak m\) 且 \(sx=0\)，与 \(I\subseteq\mathfrak m\) 矛盾。故全部局部化为零时 \(M=0\)，这一步不需要有限生成。若进一步知道 \(M\) 有限生成且 \(M(\mathfrak m)=0\)，就有 \(M_{\mathfrak m}=\mathfrak mR_{\mathfrak m}M_{\mathfrak m}\)。对局部环 \(R_{\mathfrak m}\)，Jacobson 根是其唯一极大理想；\cref{b2:thm:noncommutative-nakayama}因而使有限生成模 \(M_{\mathfrak m}\) 为零。这里从纤维消失回到局部化消失，才用到有限生成。反向都由取张量立即得到。

设 \(T\) 是局部环，极大理想为 \(\mathfrak n\)，而 \(L\) 为有限生成投射 \(T\) 模。我们想把剩余域上的基提升成 \(L\) 的基：提升 \(L/\mathfrak nL\) 的一组基，得到 \(f:T^r\rightarrow L\)。\cref{b2:cor:nakayama-generators}先保证这些提升生成 \(L\)，所以 \(f\) 满射。要进一步排除它们之间的关系，利用投射性写 \(T^r\cong L\oplus K\)，其中 \(K\) 作为有限自由模的直和因子仍有限生成。模去 \(\mathfrak n\) 后，\(f\) 已是同构，故 \(K/\mathfrak nK=0\)；再由 Nakayama 引理，\(K=0\)。于是 \(L\cong T^r\)。把 \(P\oplus P'\cong R^N\) 局部化可知 \(P_{\mathfrak p}\) 仍有限生成且投射，取 \(T=R_{\mathfrak p}\) 即得第二项。

对第三项，先检测满射性。\(\operatorname{coker}u\) 是有限生成模 \(Q\) 的商模，张量的右正合性使它的全部极大纤维为零。第一项说明 \(u\) 满射。要检测单射性，不能直接假定取纤维保持核；这里需要目标 \(Q\) 的投射性，使该满射分裂，得到 \(P\cong\ker u\oplus Q\)。核作为有限生成模 \(P\) 的直和因子仍有限生成；对分裂列取纤维仍分裂，纤维同构假设使核的全部纤维为零。再次应用第一项，得到核为零。反向由任何函子保持同构得到。
\end{proof}

\begin{lemma}[有限投射模与底环变换]\label{b2:lem:azumaya-projective-base-change}
设 \(R,S\) 为非零交换含幺环、\(R\to S\) 为环同态，\(P,Q\) 为有限生成投射 \(R\) 模，则 \(P,Q\) 的对偶模和张量积仍有限生成且投射，并有自然同构
\begin{align}
S\otimes_R\operatorname{Hom}_R(P,Q)
&\cong\operatorname{Hom}_S(S\otimes_RP,S\otimes_RQ),\label{b2:eq:azumaya-hom-base-change}\\
\operatorname{End}_R(P)\otimes_R\operatorname{End}_R(Q)
&\cong\operatorname{End}_R(P\otimes_RQ).\label{b2:eq:azumaya-end-tensor}
\end{align}
后一个映射把 \(f\otimes g\) 送到 \(p\otimes q\mapsto f(p)\otimes g(q)\)，并为代数同构。若 \(P\) 还忠实，则它是生成元，且 \(S\otimes_RP\) 是忠实有限生成投射 \(S\) 模。两个忠实有限生成投射模的张量积也忠实。
\end{lemma}
\begin{proof}
有限自由模直和分解取对偶、取底环变换或逐因子取张量后仍然分裂，所以所列模均为有限自由模的直和因子。\cref{b2:prop:finite-projective-hom-tensor}给出 \(\operatorname{Hom}_R(P,Q)\cong P^*\otimes_RQ\)。对偶的底环变换公式先在有限自由模上由坐标成立，再限制到直和因子，得到式\eqref{b2:eq:azumaya-hom-base-change}。式\eqref{b2:eq:azumaya-end-tensor}在极大理想处化为两个有限自由模的矩阵单位对应，故由\cref{b2:lem:azumaya-local-toolkit}的同构判据成立；映射的公式直接保留乘法。

现在设 \(P\) 忠实。若某个 \(P_{\mathfrak m}=0\)，取 \(P\) 的有限生成组，分别用不在 \(\mathfrak m\) 中的元素零化它们，再取这些元素的乘积，得到 \(s\notin\mathfrak m\) 且 \(sP=0\)。忠实性使 \(s=0\)，矛盾。因此每个 \(P_{\mathfrak m}\) 都自由且秩为正。

回忆\cref{b2:prop:trace-generator-criterion}中的迹理想
\[
I=\operatorname{tr}_R(P)=\sum_{\lambda\in P^*}\lambda(P)\subseteq R,
\qquad P^*=\operatorname{Hom}_R(P,R).
\]
它由全部 \(R\) 线性泛函在 \(P\) 上的取值生成。式\eqref{b2:eq:azumaya-hom-base-change}使取局部对偶与取局部化相容，故 \(I_{\mathfrak m}\) 由 \(P_{\mathfrak m}\) 的对偶泛函在模元素上的取值生成。由于局部自由秩为正，选一个基向量及其坐标泛函，便有 \(I_{\mathfrak m}=R_{\mathfrak m}\)。因此 \((R/I)_{\mathfrak m}=0\) 对全部 \(\mathfrak m\) 成立，局部消失判据给出 \(I=R\)。虽然定义迹理想时允许所有泛函，表示其中的单位元只用有限多个取值；将系数吸收到泛函中，可写
\begin{equation}\label{b2:eq:azumaya-trace-unit}
1=\sum_{i=1}^t\lambda_i(p_i).
\end{equation}
由\cref{b2:prop:trace-generator-criterion}，这就说明 \(P\) 为生成元。

底环变换后式\eqref{b2:eq:azumaya-trace-unit}仍成立。若 \(s\in S\) 零化 \(S\otimes_RP\)，对这个等式乘以 \(s\)，再利用各扩张泛函的 \(S\) 线性性，得到 \(s=0\)。因此底环变换后仍忠实。若 \(Q\) 也忠实，将两者表示单位元的等式相乘，便得到 \(P\otimes_RQ\) 上相应的等式；同一论证证明张量积忠实。
\end{proof}

忠实性在这里与有限投射性一起使用：忠实性与有限生成共同保证每个局部化非零，投射性再使局部对偶中有坐标泛函，从而得到式\eqref{b2:eq:azumaya-trace-unit}。单有忠实性既不能推出模为生成元，也不能保证任意底环变换后仍忠实；\(\mathbb Q\) 作为 \(\mathbb Z\) 模的例子见\subsecref{b2:subsec:1601:01}及下文\cref{b2:prop:C-fiber-hypothesis-boundaries}。

\ProseParagraphBreak
局部环上的自由基还可以延伸到点的一个邻域。为说明这里的邻域，给前面的素理想集合加上拓扑：对 \(f\in R\)，记 \(D(f)=\{\mathfrak p\in\operatorname{Spec}R\mid f\notin\mathfrak p\}\)。这些集合满足 \(D(1)=\operatorname{Spec}R\)、\(D(f)\cap D(g)=D(fg)\)；以它们为开集的一组基，得到素理想谱上的\term[Zariski-tuopu]{Zariski 拓扑}[Zariski topology]。集合 \(D(f)\) 称为主开集，在其中允许 \(f\) 可逆，对应于底环局部化 \(R_f=R[1/f]\)。
\symidx{\(D(f)\)：素理想谱的主开集}

\begin{proposition}[局部基的延伸]\label{b2:prop:projective-basic-open-basis}
设 \(R\) 为非零交换含幺环、\(P\) 为有限生成投射 \(R\) 模。记 \(\operatorname{Spec}R\) 为 \(R\) 的素理想集合；对 \(f\in R\)，记 \(D(f)=\{\mathfrak p\in\operatorname{Spec}R\mid f\notin\mathfrak p\}\)、\(R_f=R[1/f]\) 和 \(P_f=R_f\otimes_RP\)。每个 \(\mathfrak p\in\operatorname{Spec}R\) 都有一个主开邻域 \(D(f)\)，使 \(P_f\) 为有限自由 \(R_f\) 模。这样的主开集可以选有限多个覆盖 \(\operatorname{Spec}R\)。
\end{proposition}
\begin{proof}
\cref{b2:lem:azumaya-local-toolkit}说明 \(P_{\mathfrak p}\) 自由。其有限个基向量是 \(P\) 中元素的分式，把每个基向量乘以一个可逆分母后，仍得到一组基，故可选 \(p_1,\ldots,p_r\in P\)，使映射 \(\alpha:R^r\rightarrow P\) 在 \(\mathfrak p\) 处为同构。

由式\eqref{b2:eq:azumaya-hom-base-change}，局部逆映射可写成某个 \(\beta\in\operatorname{Hom}_R(P,R^r)\) 除以 \(s\notin\mathfrak p\)。要使它在一个主开集上已经成为逆映射，只需同时消去两个复合与恒等映射之间的误差。在 \(R_s\) 上，考察 \((\beta/s)\alpha-1\) 与 \(\alpha(\beta/s)-1\)，它们在 \(\mathfrak p\) 处为零。分别取 \(R_s^r\) 和 \(P_s\) 的有限生成组，对这有限多个值应用分式零判据，再把所得分母相乘，可以选 \(t\notin\mathfrak p\)，使两映射在 \(R_{st}\) 上都为零。因此 \(f=st\) 满足要求。有限基、局部逆映射和有限多个误差值，正是这里能够一次清去分母的有限数据。

若所有这些 \(f\) 生成的理想为真理想，取包含它的极大理想，就得到一个不在任何 \(D(f)\) 中的素理想，矛盾。因此有有限个 \(f_1,\ldots,f_u\) 及 \(a_i\in R\)，使 \(\sum_i a_if_i=1\)。任何素理想不能同时包含这些 \(f_i\)，故对应的有限个主开集已经覆盖全谱。
\end{proof}

最后这个单位理想论证对任意主开覆盖都成立。因为主开集构成拓扑的一组基，它也说明 \(\operatorname{Spec}R\) 是拟紧空间。

\subsection{纤维判据与底环变换}
\label{b2:subsec:C02:02}

设 \(R\) 为非零交换含幺环、\(A\) 为结合含幺 \(R\) 代数，且结构同态取值于中心。对 \(\mathfrak p\in\operatorname{Spec}R\)，局部化与取纤维分两步进行：
\[
\begin{tikzcd}[column sep=3.5em,row sep=2em]
R\arrow[r]\arrow[d]&R_{\mathfrak p}\arrow[r,two heads]\arrow[d]&\kappa(\mathfrak p)\arrow[d]\\
A\arrow[r]&A_{\mathfrak p}\arrow[r,two heads]&A(\mathfrak p)
\end{tikzcd}
\]
这里 \(A_{\mathfrak p}=R_{\mathfrak p}\otimes_RA\) 为局部化代数，\(A(\mathfrak p)=\kappa(\mathfrak p)\otimes_RA\) 为纤维代数；竖箭头为结构同态，两个方块均交换。第一步使 \(R\setminus\mathfrak p\) 中的元素可逆，第二步再模去 \(\mathfrak pR_{\mathfrak p}\)，因而
\[
A(\mathfrak p)\cong A_{\mathfrak p}/\mathfrak pA_{\mathfrak p}.
\]
局部化代数的底环仍是局部环 \(R_{\mathfrak p}\)，纤维代数的底环则是剩余域 \(\kappa(\mathfrak p)\)。纤维由代数自身作为模的商构造得到；即使代数嵌在矩阵代数中，也不能未经检验就把它等同于逐项代入剩余域所得的矩阵像，具体区别见\cref{b2:prop:C-generic-fiber-insufficient}。

\begin{theorem}[Azumaya 代数的纤维判据]\label{b2:thm:azumaya-fiber-criterion}
设 \(R\) 为非零交换含幺环、\(A\) 为结合含幺 \(R\) 代数，其结构同态取值于中心，且 \(A\) 作为 \(R\) 模忠实、有限生成且投射。对素理想 \(\mathfrak p\)，记 \(R_{\mathfrak p}=(R\setminus\mathfrak p)^{-1}R\)、\(A_{\mathfrak p}=R_{\mathfrak p}\otimes_RA\)、\(\kappa(\mathfrak p)=R_{\mathfrak p}/\mathfrak pR_{\mathfrak p}\)、\(A(\mathfrak p)=\kappa(\mathfrak p)\otimes_RA\)。以下条件等价：
\begin{enumerate}
\item \(A\) 为 Azumaya \(R\) 代数；
\item 对每个极大理想 \(\mathfrak m\)，\(A(\mathfrak m)\) 为中心单 \(\kappa(\mathfrak m)\) 代数；
\item 对每个素理想 \(\mathfrak p\)，\(A_{\mathfrak p}\) 为 Azumaya \(R_{\mathfrak p}\) 代数。
\end{enumerate}
此外，若 \(A\) 满足这些条件，且 \(R\to S\) 为非零交换含幺环的同态，则 \(S\otimes_RA\) 为 Azumaya \(S\) 代数。
\end{theorem}
\begin{proof}
Azumaya 条件要检验的是 \(\mu_A\) 是否同构，而纤维上的中心单条件正好检验它的底环变换。\cref{b2:lem:azumaya-projective-base-change}保证 \(A\otimes_RA^{\mathrm{op}}\) 与 \(\operatorname{End}_R(A)\) 都有限生成且投射，并使 \(\mu_A\) 在 \(\kappa(\mathfrak m)\) 上的底环变换恰为 \(\mu_{A(\mathfrak m)}\)。该引理的忠实性论证还给出 \(A_{\mathfrak m}\) 的自由秩为正，所以 \(A(\mathfrak m)\ne0\)；仅有投射性并不能保证这一点。由\cref{b2:prop:azumaya-field-matrix}，该纤维映射同构等价于纤维代数中心单；再由\cref{b2:lem:azumaya-local-toolkit}检测同构，第一、二项等价。

若第一项成立，对任意 \(R\rightarrow S\) 取底环变换。\cref{b2:lem:azumaya-projective-base-change}给出忠实有限投射性，又有
\[
(S\otimes_RA)\otimes_S(S\otimes_RA)^{\mathrm{op}}
\cong S\otimes_R(A\otimes_RA^{\mathrm{op}}),
\]
以及自同态代数的底环变换同构。在这些识别下，\(\mu_{S\otimes_RA}\) 是 \(S\otimes_R\mu_A\)，故为同构。这里不要求 \(S\) 平坦：任何函子都保持同构，而有限投射性保证所用的自同态变基公式，迹理想中的单位元等式保证变基后的忠实性。这证明底环变换断言，也证明第一项蕴含第三项。若第三项成立，再对 \(R_{\mathfrak m}\rightarrow\kappa(\mathfrak m)\) 作底环变换便得第二项，等价性证毕。
\end{proof}

判据要求检验所有极大纤维；只在分式域上或某个主开集上成为矩阵代数，仍不足以推出原代数 Azumaya，见\cref{b2:prop:C-generic-fiber-insufficient}。局部模的同构判据还可检测有限投射代数中单个元素的可逆性，见\cref{b2:prop:C-unit-fiber-criterion}。

\begin{corollary}\label{b2:cor:azumaya-endomorphism}
设 \(R\) 为非零交换含幺环。若 \(P\) 是忠实有限生成投射 \(R\) 模，则 \(\operatorname{End}_R(P)\) 是 Azumaya \(R\) 代数。
\end{corollary}
\begin{proof}
由\cref{b2:lem:azumaya-projective-base-change}，自同态模为有限投射模。若 \(r\) 零化所有自同态，它就零化恒等映射，故零化 \(P\)，于是 \(r=0\)。每个极大纤维由式\eqref{b2:eq:azumaya-hom-base-change}同构于非零有限维向量空间的自同态代数，所以是中心单代数。纤维判据给出结论。
\end{proof}

\begin{proposition}\label{b2:prop:azumaya-center}
设 \(R\) 为非零交换含幺环、\(A\) 为 Azumaya \(R\) 代数，则 \(A\) 的中心为 \(R1_A\)。此外，对任意忠实有限生成投射 \(R\) 模 \(P\)，有
\[
Z\bigl(\operatorname{End}_R(P)\bigr)
=\{r\operatorname{id}_P\mid r\in R\}.
\]
这两处的标量表示都唯一。
\end{proposition}
\begin{proof}
先设 \(P\) 为忠实有限生成投射 \(R\) 模，\(T\in\operatorname{End}_R(P)\) 与全部自同态交换。取式\eqref{b2:eq:azumaya-trace-unit}中的 \(p_i,\lambda_i\)。要证明 \(T(x)\) 是同一个标量乘以 \(x\)，可先乘上这些取值后再求和；为把 \(T\) 从 \(x\) 移到 \(p_i\)，让它与自同态 \(y\mapsto\lambda_i(y)x\) 交换，再在 \(p_i\) 处取值，得到
\[
\lambda_i(p_i)T(x)=\lambda_i\bigl(T(p_i)\bigr)x.
\]
求和说明 \(T(x)=rx\)，其中 \(r=\sum_i\lambda_i\bigl(T(p_i)\bigr)\in R\) 不依赖于 \(x\)。反过来，标量乘法与每个 \(R\) 线性自同态交换；忠实性保证不同标量给出不同自同态。由此得到关于自同态代数中心的断言。

若 \(z\in Z(A)\)，左乘 \(z\) 与所有左右乘算子交换。\(\mu_A\) 满射，因而它与全部 \(R\) 线性自同态交换。取 \(P=A\)，由刚证结论有 \(L_z=r\operatorname{id}_A\)，在 \(1_A\) 处取值得 \(z=r1_A\)。反向包含是 \(R\) 代数的约定，而忠实性保证这种标量表示唯一。
\end{proof}

纤维判据中的“中心单”不能换成“在剩余域上已经分裂”而仍声称二者等价。对基环 \(\mathbb R\)，唯一素理想为零，局部化和剩余域都不改变基环；\(\mathbb H\) 的纤维仍为除代数。换到 \(\mathbb C\) 后它才成为 \(M_2(\mathbb C)\)。因此 Azumaya 性并不意味着在 Zariski 局部化后就能选出矩阵单位。

\ProseParagraphBreak
这个例子也提示了分裂所需的另一种局部操作：除了使指定元素可逆，还允许剩余域上的有限可分扩张。平展代数同时包含这两种操作。本附录用以下仿射条件描述它：设 \(R\to S\) 为交换含幺环的同态，若 \(S\) 作为交换 \(R\) 代数有有限个生成元和有限条多项式关系，作为 \(R\) 模平坦，且对每个素理想 \(\mathfrak p\)，\(\kappa(\mathfrak p)\otimes_RS\) 是有限个有限可分 \(\kappa(\mathfrak p)\) 域扩张的直积，则称 \(S\) 为\term[pingzhan-daishu]{平展代数}[étale algebra]。域上的有限可分扩张与底环的主开局部化都满足这些条件。

\ProseParagraphBreak
这里允许零纤维，即空直积。例如 \(\kappa(\mathfrak p)\otimes_RR_f\cong\kappa(\mathfrak p)_f\)，若 \(f\in\mathfrak p\)，它要求把剩余域中的零变成单位，因而为零环。这表示 \(\operatorname{Spec}R_f\) 没有点映到 \(\mathfrak p\)，即这个主开集没有覆盖该点。因此下面除了要求各处平展分裂，还必须要求这些谱映射覆盖全谱。

\begin{theorem}[平展局部分裂]\label{b2:thm:azumaya-etale-splitting}
设 \(R\) 为非零交换含幺环、\(A\) 为结构同态取值于中心的结合含幺 \(R\) 代数，且 \(A\) 作为 \(R\) 模忠实、有限生成且投射。则 \(A\) 是 Azumaya \(R\) 代数，当且仅当存在有限族非零交换平展 \(R\) 代数 \(S_1,\ldots,S_t\)，满足以下两项条件：
\begin{enumerate}
\item 每个素理想 \(\mathfrak p\in\operatorname{Spec}R\) 都是某个 \(S_i\) 的素理想在 \(R\) 中的收缩，即这些谱映射的像联合覆盖 \(\operatorname{Spec}R\)。
\item 对每个 \(i\)，存在正整数 \(n_i\) 和 \(S_i\) 代数同构
\[
S_i\otimes_RA\cong M_{n_i}(S_i).
\]
\end{enumerate}
\end{theorem}
\begin{proofsketch}
正向需要把纤维上的可分分裂代数提升到点附近，这一提升不由前面的局部模论直接给出。这里只说明所调用的外部结论及其用途：每个点都有开邻域 \(U\) 以及有限平展满射 \(U'\to U\)，使 \(A\) 在 \(U'\) 上成为矩阵代数。这一提升及分裂构造见\cite[Théorème 5.1, Proposition 5.4, Théorème 5.7]{Grothendieck1966BrauerI}；此处省去提升步骤的证明。把 \(U\) 缩小为主开集，\(U'\) 仍为仿射，从而得到所需的平展 \(R\) 代数。\(\operatorname{Spec}R\) 的拟紧性使覆盖可以取为有限族。

反向可以用前面的模论完整证明。令 \(S=\prod_i S_i\)，各 \(S_i\) 平坦，有限直积 \(S\) 也平坦。谱覆盖还保证它的张量函子检测非零模。确实，对 \(0\ne x\in M\)，令 \(I=\operatorname{Ann}_R(x)\)，并选包含 \(I\) 的极大理想 \(\mathfrak m\)。覆盖条件给出某个 \(S_j\) 的素理想 \(\mathfrak q\)，收缩为 \(\mathfrak m\)；于是 \(IS_j\subseteq\mathfrak q\)，故 \(S_j\otimes_RRx\cong S_j/IS_j\ne0\)。平坦性使这个模嵌入 \(S_j\otimes_RM\)，所以 \(S\otimes_RM\ne0\)。因此 \(S\) 不仅平坦，而且张量后一个模为零会迫使原模为零，这就是这里所需的忠实平坦性。

在每个 \(S_i\) 上，\(\mu_A\) 经底环变换成为矩阵代数的左右乘同构。由式\eqref{b2:eq:azumaya-hom-base-change}，\(S\otimes_R\mu_A\) 因而为同构。平坦性使取张量保持核、余核，所以
\[
S\otimes_R\ker\mu_A=0,
\qquad S\otimes_R\operatorname{coker}\mu_A=0.
\]
忠实平坦性使 \(\ker\mu_A=0\) 且 \(\operatorname{coker}\mu_A=0\)。因此 \(\mu_A\) 同构，得到 Azumaya 性。
\end{proofsketch}

有限族指覆盖中代数的个数有限，不要求每个 \(S_i\) 都是有限 \(R\) 模。矩阵大小也可随基谱的不同部分变化。平展局部分裂是 Azumaya 代数的几何刻画；它允许像 \(\mathbb R\to\mathbb C\) 这样的可分扩张，因而比只取主开集的 Zariski 局部分裂更一般。

\begin{proposition}[纤维检验的假设边界]\label{b2:prop:C-fiber-hypothesis-boundaries}
非零 \(\mathbb Z\) 模 \(\mathbb Q\) 忠实，且对每个素数 \(p\)，有
\[
\mathbb Q_{(p)}\cong\mathbb Q\ne0,
\qquad
\mathbb F_p\otimes_{\mathbb Z}\mathbb Q=0.
\]
因而不能对任意模用全部极大纤维的消失检测原模的消失，也不能对任意忠实模断言底环变换后仍忠实。另设 \(k\) 为任意域、\(R=k[\varepsilon]/(\varepsilon^2)\)，则商映射 \(\rho:R\to k=R/(\varepsilon)\) 的唯一极大纤维为同构，但 \(\ker\rho=(\varepsilon)\ne0\)，且目标 \(k\) 不是投射 \(R\) 模。
\end{proposition}
\begin{proof}
非零整数不能零化 \(\mathbb Q\)，所以它作为 \(\mathbb Z\) 模忠实。任意不在 \((p)\) 中的整数在 \(\mathbb Q\) 上已经可逆，局部化的映射 \(q/s\mapsto q/s\) 给出 \(\mathbb Q_{(p)}\cong\mathbb Q\)。另一方面，\(p\mathbb Q=\mathbb Q\)，故
\[
\mathbb F_p\otimes_{\mathbb Z}\mathbb Q
\cong\mathbb Q/p\mathbb Q=0.
\]
这不违反\cref{b2:lem:azumaya-local-toolkit}，因为 \(\mathbb Q\) 并非有限生成 \(\mathbb Z\) 模：有限多个有理数的分母有公共倍数 \(d\)，它们的整数线性组合都在 \(d^{-1}\mathbb Z\) 中，不能包含 \(1/(2d)\)。所得零模在非零环 \(\mathbb F_p\) 上也不忠实。

对双数环 \(R\)，元素 \(a+b\varepsilon\) 在 \(a\ne0\) 时可逆，其逆为 \(a^{-1}-a^{-2}b\varepsilon\)，而 \((\varepsilon)\) 的商为域 \(k\)，所以 \((\varepsilon)\) 是唯一极大理想。对 \(\rho\) 在这个点取纤维，得到 \(k\otimes_RR\to k\otimes_Rk\)，即 \(k\) 上的恒等映射；原映射却有非零核 \((\varepsilon)\)。若 \(k\) 投射，\(\rho\) 就有 \(R\) 线性截面 \(s:k\to R\)。从 \(\rho(s(1))=1\) 可写 \(s(1)=1+b\varepsilon\)，但 \(k\) 被 \(\varepsilon\) 零化，线性性要求
\[
0=s(\varepsilon\cdot1)=\varepsilon s(1)=\varepsilon,
\]
矛盾。因此，纤维上的同构没有排除原核，正是目标投射性的缺失所允许的现象。
\end{proof}

\begin{proposition}[泛纤维不足以检验 Azumaya 性]\label{b2:prop:C-generic-fiber-insufficient}
设 \(k\) 为任意域、\(R=k[t]\)，并令
\[
A=\left\{
\begin{pmatrix}a&b\\tc&d\end{pmatrix}
\ \middle|\ a,b,c,d\in R
\right\}\subseteq M_2(R).
\]
代数 \(A\) 含幺，底模忠实且自由，秩为 \(4\)。它满足 \(A[1/t]\cong M_2(R[1/t])\)，从而泛纤维 \(k(t)\otimes_RA\cong M_2(k(t))\)；但在极大理想 \((t)\) 处的纤维 \(A/tA\) 不是单代数，所以 \(A\) 不是 Azumaya \(R\) 代数。将矩阵条目在 \(t=0\) 处求值的映射在该纤维上并不单射。
\end{proposition}
\begin{proof}
矩阵乘法保持左下条目为 \(t\) 的倍数，且 \(1_A=E_{11}+E_{22}\)，所以 \(A\) 确为含幺子代数。令
\[
e=E_{11},\qquad f=E_{22},\qquad x=E_{12},\qquad y=tE_{21}.
\]
每个元素唯一写成 \(ae+df+bx+cy\)，故这四个元素为 \(R\) 基，底模自由且忠实。使 \(t\) 可逆后，\(t^{-1}y=E_{21}\)，四个矩阵单位均在局部化中，因而 \(A[1/t]=M_2(R[1/t])\)。再把底环扩张到 \(k(t)\) 即得泛纤维的断言。

由于 \(A\) 的这四个元素构成自由基，它们的商类 \(\bar e,\bar f,\bar x,\bar y\) 构成 \(A/tA\) 的 \(k\) 基，特别是 \(\bar y\ne0\)。在 \(A\) 中有
\[
x^2=y^2=0,\qquad xy=te,\qquad yx=tf.
\]
又有 \(ex=x=xf\)、\(fy=y=ye\)，其余用 \(e,f\) 从两侧乘 \(x,y\) 的结果为零。因此在 \(A/tA\) 中，\(J=k\bar x+k\bar y\) 是非零真双边理想，且 \(J^2=0\)。纤维不是单代数，\cref{b2:thm:azumaya-fiber-criterion}便排除了 \(A\) 的 Azumaya 性。

把 \(A\) 中的矩阵条目在 \(t=0\) 处求值，得到满射到上三角矩阵代数的映射。它零化 \(tA\)，所以诱导 \(A/tA\) 上的映射，却把 \(\bar y\) 送到零。内在纤维仍保留基向量 \(y\) 的非零商类；外在的条目求值丢掉了它，故这两个构造不能混同。
\end{proof}

\begin{proposition}[有限投射代数中的单位检验]\label{b2:prop:C-unit-fiber-criterion}
设 \(R\) 为非零交换含幺环、\(A\) 为结构同态取值于中心的结合含幺 \(R\) 代数，且 \(A\) 作为 \(R\) 模有限生成且投射。对 \(a\in A\)，\(a\) 是 \(A\) 的单位，当且仅当对每个极大理想 \(\mathfrak m\)，它在纤维代数 \(A(\mathfrak m)=(R/\mathfrak m)\otimes_RA\) 中的像都是单位。这里不要求 \(A\) 忠实或满足 Azumaya 条件。
\end{proposition}
\begin{proof}
若 \(ab=ba=1_A\)，取任一纤维后仍有这两个等式，故正向成立。反之，左乘映射 \(L_a:A\to A\) 是 \(R\) 线性映射；在各极大纤维中，它变成相应单位的左乘映射，因而为同构。\cref{b2:lem:azumaya-local-toolkit}使 \(L_a\) 本身为同构。令 \(b=L_a^{-1}(1_A)\)，则 \(ab=1_A\)。为得到另一侧的逆，计算
\[
L_a(ba-1_A)=a(ba-1_A)=(ab)a-a=0.
\]
由 \(L_a\) 的单射性，\(ba=1_A\)，所以 \(a\) 可逆。
\end{proof}
